% Gerado de tcc.md (exportado do Claude Docs) por ferramentas/gera_latex.php em 12/11/2026 21:45.
% Para mudar o texto, edite este arquivo; se gerar de novo a partir do documento, as edições feitas aqui se perdem.

\noindent O Plano Educacional Individualizado (PEI) organiza metas de aprendizagem para alunos público-alvo da educação especial, mas o acompanhamento dessas metas costuma depender de registros escritos pelo professor, sem dados que mostrem a evolução ao longo do tempo. Este trabalho apresenta o UNIA (Unir Necessidades Individuais e Aprendizagem), um sistema web que apoia o professor no acompanhamento das metas do PEI de crianças dos anos iniciais e mede o desempenho delas por meio de mini games. A proposta parte do AppPEI, recurso educacional digital em planilha eletrônica desenvolvido na Universidade do Estado do Rio de Janeiro, do qual foram aproveitadas as quatro áreas da matriz de habilidades e os três níveis de suporte: não realiza, realiza com ajuda e realiza sem necessidade de suporte. A pesquisa é aplicada e exploratória, com revisão bibliográfica e desenvolvimento incremental de software em PHP, no padrão Model-View-Controller (MVC), com banco de dados MySQL, testado no MariaDB 10.4 que acompanha o XAMPP. No UNIA, o professor cadastra o aluno e suas metas, e o aluno responde a sessões de cinco rodadas em 12 formatos de mini game, distribuídos em 24 atividades, seis por área. O sistema confere cada resposta e registra acerto, uso de dica e tempo; ao final, classifica o nível de suporte da sessão e calcula, por meta, o acerto médio, a proporção de sessões sem suporte, a proporção de sessões com dica e a variação entre a primeira e a segunda metade das sessões, apresentados em tabelas e gráficos acessíveis. A verificação técnica incluiu 63 testes automatizados das regras, todos aprovados, e um teste de ponta a ponta em banco de dados temporário; uma avaliação heurística pelos princípios de Nielsen apontou dez problemas de usabilidade, dois deles graves, que foram corrigidos. O sistema ainda não foi avaliado com professores e alunos; a regra de classificação e os parâmetros adotados são hipóteses do protótipo, cuja validação pedagógica fica como próximo passo. Conclui-se que é viável integrar a medição automática do desempenho ao acompanhamento do PEI, transformando registros pontuais em séries que mostram a evolução de cada meta.

\vspace{\baselineskip}
\noindent\textbf{Palavras-chave:} Plano Educacional Individualizado. Educação especial. Jogos digitais. Acompanhamento da aprendizagem. Sistema web.
